\section{Implementation}
\label{sec:implementation}
This section describes how the design of \cref{sec:concepts} is realised in code. It concentrates on the
decisions that shaped the implementation rather than on the code itself; the full source is available in the
project repository\footnote{\url{https://github.com/FHNW-Security-Lab-Student-Projects/ip6-autonome-netzwerkadministration}}.

\subsection{Agent Framework and Harness}
\label{sec:impl-framework}
The agent harness introduced conceptually in \cref{sec:agent-loop} is provided in this work by
\emph{Pydantic AI}, a Python agent framework \parencite{pydantic_ai}. Pydantic AI implements the agent loop,
drives the function-calling exchange with the model, manages the message history across iterations and
exposes the tools of each agent.

Every agent of \cref{sec:agent-roles} is one Pydantic AI agent, defined by four things: a name, an
instruction text (its system prompt), a model and a set of tools. The entire multi-agent system runs in a
single Python process. Agents are plain objects that call each other directly, there is no network protocol
or separate service per agent. This keeps the system simple to start, debug and reason about, an earlier
design based on an inter-agent network protocol was dropped for exactly this reason.

\subsection{Model Access}
\label{sec:impl-models}
All models are accessed through \emph{OpenRouter}, a gateway that exposes models from different vendors
behind one API \parencite{openrouter}. Switching the model an agent uses therefore only means changing a
model identifier, no vendor-specific code is involved. This is what makes the experiment matrix of
\cref{sec:experiment-matrix} practical: the same system can be swept across six models from three price
tiers without touching the agent logic.

The selected model is a runtime parameter. The orchestrator passes the active model down into every
sub-agent run, so a whole troubleshooting session runs on one model, which is required to attribute an
outcome to a single tier (\cref{sec:model-tier-concept}).

Two settings are applied centrally to every agent. First, the reasoning effort is fixed to the same level
for all models, so that model-versus-model comparisons are not distorted by each provider picking its own
default reasoning depth. Second, usage accounting is enabled on every request, so each model response
carries its native token counts. These counts are the basis for the cost and token metrics of
\cref{sec:metrics}.

\subsection{Delegation as Tool Calls}
\label{sec:impl-delegation}
Delegation between the orchestrator and its sub-agents (\cref{sec:architecture}) is realised as ordinary
function calling: each sub-agent is registered on the orchestrator as a tool. The tool description tells
the orchestrator's model what the sub-agent can do and when to use it. When the model calls the tool, the
wrapper runs the sub-agent's own agent loop to completion and returns its final answer as the tool result.

Sub-agents are stateless: they never see the orchestrator's conversation and keep no memory between calls.
The orchestrator's instructions therefore require every delegated request to be self-contained, carrying
the user's goal, the relevant details and any findings from earlier sub-agent calls. Read-only sub-agents
may be called in parallel when a request involves independent lookups.

One agent role is deliberately not backed by an LLM. The topology agent is implemented as a background task
that queries the LLDP neighbours of every device once a minute, compares them against the intended topology
definition of the lab and caches the result. The orchestrator's topology tool simply returns this cache.
Since topology retrieval is deterministic, a model in this path would add cost and latency without adding
any capability.

\subsection{Tool Servers and Device Access}
\label{sec:impl-tools}
Following the design of \cref{sec:tool-layer}, each tool-using agent is served by its own MCP server
(\cref{sec:mcp}), implemented with \emph{FastMCP} \parencite{fastmcp} and started as a subprocess of the
main process. Tool names are prefixed per domain, so tools of different agents cannot be confused.

All interaction with the lab devices goes through the JSON-RPC interface of Nokia SR Linux
\parencite{srlinux_jsonrpc}, wrapped in a single shared client module. The tools offer two routes to device
data: executing formatted \texttt{show} commands, and reading a specific YANG path directly from the state
or running-configuration datastore. The agent instructions explain when each route is appropriate. To keep
the models from inventing plausible-looking but invalid commands, a reference of valid SR Linux read
commands is embedded in the network agent's system prompt. When a device still rejects a command, the
failure is detected in the tool output and recorded, this is the source of the invalid-command count of
\cref{sec:metrics}.

For syslog, all devices log to a central \emph{Loki} instance \parencite{grafana_loki}, which the syslog
agent queries through its MCP tool. A practical problem here is time: an LLM has no reliable clock and its
prompt may be minutes old by the time a tool runs. The tools therefore accept time expressions such as
\enquote{now}, a relative offset (\enquote{15m ago}) or an absolute timestamp, and resolve them against the
real clock at call time. The model only ever states its intent and never computes timestamps itself.

\subsection{Context Management}
\label{sec:impl-context}
Device output is the dominant source of token growth: every iteration of an agent loop re-sends the full
message history, so large tool returns are paid for again on every subsequent step. Observed runs of the
network agent reached several hundred thousand input tokens for a single query before countermeasures were
in place. Two measures address this, both following from \cref{sec:tokens-context-cost}.

The first works at the source. Tool responses are serialised compactly and cut off at a fixed size; a
truncated response ends with a note telling the model to re-query a narrower path instead. The instructions
additionally steer the model towards narrow queries (a single leaf across all interfaces rather than the
whole interface tree).

The second works on the history. Before each model request, a processor estimates the current prompt size,
anchored on the token count the provider reported for the previous request. Only when the estimate
approaches the model's context window (whose size is looked up per model, since the swept models differ
considerably) does it act: the oldest oversized tool returns are replaced by a one-line stub stating which
tool was called and that it can be re-invoked if the output is needed again. The most recent returns and
all error messages are always kept verbatim, errors carry the correction the model needs to avoid repeating
a bad call. This keeps long runs inside the context window, reduces cost and counters the
\enquote{lost in the middle} degradation without hiding any information the model cannot get back.

\subsection{Guarding Configuration Changes}
\label{sec:impl-config}
The configuration agent is the only write path into the network, and its use is guarded by a two-step
workflow encoded in the orchestrator's instructions. A change request is first delegated as
validation-only: the agent loads the commands into the device's candidate configuration and returns the
resulting diff without committing. The orchestrator shows this diff to the user, and only after explicit
approval does it send the apply request, which must carry the exact commands from the validation step. A
request without a clear prefix is treated as validation, so the default is always the safe path. The
autonomous investigation orchestrator has no access to the configuration agent at all
(\cref{sec:agent-roles}).

\subsection{Autonomous Investigations and Handover}
\label{sec:impl-investigations}
The syslog use case of \cref{sec:autonomous-syslog-usecase} is driven by a background task that polls Loki
every 30 seconds for messages of severity error and above. A matching event opens an investigation, keyed
by device and message so that a repeating log line does not spawn duplicates, and starts a triage run of
the investigation orchestrator in the background. Because this orchestrator lacks the tools to open or
continue investigations, an investigation can never trigger further investigations recursively.

Each investigation is persisted to disk as a record containing an identifier, the affected device, the
triggering log line and its timestamp, the current status, a summary, and, crucially, the complete
serialised message history of the agent run. This history is what realises the handover described in
\cref{sec:autonomous-syslog-usecase}: the user-facing orchestrator has tools to list, inspect and continue
investigations, and continuing a case loads the stored history and resumes the investigation orchestrator
on top of it. The model therefore still holds all earlier findings and tool results and continues from
where the autonomous run stopped, instead of re-running the same commands.

\subsection{Interfaces and Instrumentation}
\label{sec:impl-interfaces}
The system has three entry points: a terminal chat for development, a web chat interface with a model
selector (built on the framework's bundled web UI), and a status page showing all open investigations. In
the interactive interfaces the orchestrator keeps its message history across turns, so a troubleshooting
session is one continuous conversation.

All agent runs, tool calls and device requests are traced with \emph{Logfire} \parencite{logfire}. The
traces show, per run, which sub-agents were called with which requests and what each tool returned, which
was the main debugging instrument during development and was also used to verify the recorded token and
cost figures.

The experiments of \cref{sec:methodology} reuse the interactive system unchanged, except that the
configuration agent and the investigation tools are disabled through feature flags, so the measured system
is exactly the read-only troubleshooting core. The delegation wrappers record the usage of every sub-agent
run into the experiment log, from which the metrics of \cref{sec:metrics} are derived.
